\chapter{Un siglo de defensas}

\label{cap:defensas}
El río La Caldera baja de la cornisa, se bifurca en tres brazos cuatro kilómetros aguas
arriba del pueblo y vuelve a juntarse después. Cada tanto se lleva algo. Este capítulo sigue
lo que el Estado hizo al respecto entre 1910 y 2026, y llega a una conclusión incómoda: en
ciento dieciséis años la respuesta fue siempre la misma, y el propio Estado dejó escrito tres
veces que no alcanzaba.

Va antes que los tres capítulos siguientes porque los explica. La expropiación, la línea de ribera y el amparo son episodios de esta serie, no su comienzo.

\section{Enero de 1910}

La primera pieza de la historia del río está en el Boletín de 1910, y dice casi todo.

\begin{ficha}{Decreto del 21 de enero de 1910}
«Vista la solicitud presentada por \textbf{los vecinos del pueblo de La Caldera}, haciendo
presente la urgente necesidad que hay de defender la población, capital del departamento,
y el camino que conduce de esta ciudad a ese pueblo, \textbf{de las invasiones del Río del
mismo nombre}, y estando comprobado que realmente existe el \textbf{peligro inmediato,
puesto de manifiesto por las últimas crecientes}»\ldots

\textbf{Art.\ 1º:} destina \$2.000 «para llenar la enunciada necesidad».

\textbf{Art.\ 2º:} nombra una comisión de siete vecinos —Augusto Regis\index{Regis, Augusto}, Abelardo Figueroa\index{Figueroa, Abelardo},
Silvano Murúa\index{Murúa, Silvano Ignacio}, el doctor Rafael Usandivaras, Daniel Linares\index{Linares, Daniel}, Ángel Fernández\index{Fernández, Ángel} y Jacinto
Guerrero— para que, de acuerdo con el ingeniero inspector de Obras Públicas, «proyecte y
realice los trabajos pertinentes».

\textbf{Art.\ 3º:} encarga a esa comisión «requerir de la \textbf{Comisión Municipal} del
Departamento toda la \textbf{ayuda pecuniaria} que fuera posible».
\end{ficha}

Léase el considerando cambiando la fecha. Los vecinos piden que se defienda el pueblo de
las invasiones del río; se comprueba el peligro inmediato; las últimas crecientes lo
pusieron de manifiesto. Ciento nueve años antes del amparo colectivo, la relación entre el
pueblo de La Caldera y su río ya estaba planteada en esos términos y ya había sido llevada
al Estado provincial por iniciativa vecinal.

La estructura de la respuesta es también la misma que el capítulo \ref{cap:amparo} documenta un siglo
después: los vecinos piden, la Provincia asigna fondos y nombra una comisión, el municipio
cofinancia lo que pueda. En 2019 los vecinos fueron a la justicia en lugar de al
gobernador, y la sentencia ordenó un plan en lugar de unos trabajos, pero el reparto de
papeles no se movió.

Hay un cuarto elemento que se parece todavía más. El decreto argumenta que, aunque el
camino es nacional, la Provincia debe intervenir «cuando como en este caso el Gobierno de
la Nación no tiene votados fondos especiales a este objeto». Una jurisdicción actuando
porque otra no llega: exactamente lo que ocurre cuando el municipio construye badenes con
maquinaria provincial, y cuando el río separa dos municipios y ninguno lo administra
entero.

\subsection*{La obra se hizo, y costó más}

Dieciséis meses después, el Boletín publica la continuación. Un decreto del 16 de mayo de
1911 asigna \$958,65 adicionales «para cubrir el déficit que aparece en las cuentas
presentadas a favor del señor Daniel Linares\index{Linares, Daniel}, presidente de la comisión encargada de la
ejecución», porque los dos mil pesos originales «resultaron insuficientes».

Cuatro datos. La obra se hizo: hay cuentas
presentadas, hay quien las rinde y hay un déficit que cubrir. Costó casi la mitad más: dos
mil pesos presupuestados, \$2.958,65 finales, un sobrecosto del cuarenta y ocho por ciento
en la primera obra hidráulica documentada del departamento. Se ejecutó sin autorización
legislativa previa, y el artículo 2º la pide después —para ese decreto y también para el de
1910, que dieciséis meses más tarde seguía sin ratificar. Y el decreto de 1910 aparece aquí
fechado el 20 de enero, no el 21. La relectura de 1910 zanjó de qué lado está el error: en la
edición del 5 de febrero de 1910, leída sobre la imagen, el decreto dice «Salta, Enero 21 de
1910». \textbf{La fecha del acto es el 21; el que se equivoca, al citarlo, es el decreto de
1911.} Es la misma clase de error que este libro documenta un siglo después, cuando una
resolución cita mal el expediente de otra.

\subsection*{Quiénes eran los siete}

Los nombres de esa comisión no son siete vecinos cualesquiera. Abelardo Figueroa\index{Figueroa, Abelardo} y Augusto
Regis\index{Regis, Augusto} figuran en la nómina de 1908 de los Jurados de Reclamos sobre patentes generales del departamento —Figueroa como titular,
Regis\index{Regis, Augusto} como suplente—, y Figueroa fue designado en 1909 miembro de la Comisión Municipal en
reemplazo de Silvano Murúa\index{Murúa, Silvano Ignacio}, que renunció y que también integra esta comisión. Ángel
Fernández era el Presidente Municipal que había elevado el cuadro de rentas. Jacinto
Guerrero había sido nombrado comisario suplente de policía del departamento en febrero de
1909.

La comisión encargada de proyectar las defensas del río estaba integrada, entonces, por el
presidente municipal, un miembro de la comisión municipal que era además jurado, el que había
renunciado a ese cargo, otro jurado y el comisario suplente; completan los siete Daniel
Linares\index{Linares, Daniel}, que la presidió, y el doctor Rafael Usandivaras. En un departamento de pocos miles de habitantes, quienes atienden la crecida son los mismos que gobiernan, juzgan y vigilan.

De ahí no se sigue ninguna denuncia: en 1910 no había otra forma de hacerlo. Se sigue una
observación de estructura, que el capítulo \ref{cap:politica} retoma. La concentración de funciones en
pocas manos no es un rasgo reciente del departamento: es su forma histórica de
administrarse.

\subsection*{Y veinticinco años después, otra vez}

En enero de 1935 el ciclo se repite completo, y esta vez el decreto explica con precisión
técnica lo que estaba pasando.

\begin{ficha}{Decreto del 23 de enero de 1935 --- Expediente 181-C}
Visto el proyecto de defensas del pueblo de La Caldera elevado por la Dirección General de
Obras Públicas, y considerando:

«Que en efecto, dada la última creciente \ldots\ pone en evidencia \textbf{el peligro en que
se encuentra el pueblo La Caldera} ya que la orilla \ldots\ ha sido rellenado y parte
erosionado \textbf{quedando reducida la barranca a una pequeña altura}, amenazando el brazo
respectivo \textbf{volcarse a la calle principal} donde tendría fácil escurrimiento por la
configuración del terreno.»

«Que la solución más inmediata sería \textbf{un espigón de "pie de gallo"} a construirse con
urgencia \ldots\ tras una defensa últimamente destruida \ldots\ lo que torcería el cauce a
la orilla opuesta que era el lecho anterior.»

«Que \textbf{la solución definitiva} --- dada la duración limitada de estas defensas ---
sería la construcción de \textbf{un espigón macizo} cuya longitud, constitución y ubicación
debería estudiarse; probablemente convendría colocarlo en la primera bifurcación,
\textbf{alejaría por completo el peligro}.»

«Que el Poder Ejecutivo juzga que, para la urgencia de las defensas \textbf{y ante el
ofrecimiento de los mismos vecinos afectados}, la construcción del espigón de emergencia
podría efectuarse con un gasto aproximado de quinientos pesos \textbf{y ser llevado a cabo
por la Municipalidad de La Caldera}.»

Se liquidan \$500 a favor del presidente de la Comisión Municipal, \textbf{don Augusto
Regis\index{Regis, Augusto}}, con supervisión técnica de Obras Públicas y cargo de rendir cuenta documentada.

\textit{El impreso fecha el decreto en enero de 1934, pero el expediente es del año 1935 y el
decreto siguiente de la misma columna lleva 1935. Es un error del original.}
\end{ficha}

Este documento merece leerse dos veces, porque contiene entero el problema que el libro
describe.

\textbf{Primero, el diagnóstico es correcto y es técnico.} La orilla fue rellenada y
erosionada, la barranca perdió altura, y el brazo del río amenaza volcarse sobre la calle
principal porque el terreno se lo permite. No es una descripción vaga del peligro: es
hidráulica elemental, bien hecha, en 1935.

\textbf{Segundo, el propio decreto distingue entre la solución de emergencia y la
definitiva.} El espigón de pie de gallo es urgente y de duración limitada. La solución
definitiva sería un espigón macizo, colocado probablemente en la primera bifurcación, que
«alejaría por completo el peligro». El Estado sabía, y lo puso por escrito, que lo que
estaba autorizando no resolvía nada.

\textbf{Tercero, se hace lo urgente.} Quinientos pesos, la municipalidad como ejecutora, el
presidente como responsable de la rendición. El espigón macizo requería un estudio que el
decreto no encarga.

\textbf{Y cuarto, lo pagan en parte los vecinos.} El considerando dice «ante el ofrecimiento
de los mismos vecinos afectados»: como en 1910, la iniciativa y parte del esfuerzo vienen de
quienes viven ahí.

De modo que la serie del río queda así. En 1910 los vecinos piden y la Provincia paga dos
mil pesos. En 1911 hay que cubrir el sobrecosto. En 1935 la barranca está erosionada, se
construye un espigón de emergencia con quinientos pesos y el propio Estado consigna que la
solución definitiva es otra y que hay que estudiarla. En 2014 la Provincia gasta casi dos
millones en proteger las cámaras del acueducto. En 2018 y 2019 el río se lleva casas. En
2021 una sentencia ordena un plan que privilegie infraestructura verde sobre gris. En enero
de 2026 el municipio propone represas de retardo, y a comienzos de febrero construye bateas con maquinaria pesada.

Y el espigón de emergencia tampoco entró en presupuesto. Un mes después, el 25 de febrero
de 1935, un nuevo decreto en acuerdo de ministros amplía la partida en doscientos pesos
más, «para cancelar los gastos totales a que asciende el costo de las obras de defensas del
pueblo de La Caldera», a pedido del propio presidente de la comisión municipal y con informe
previo de Contaduría General.

De quinientos pesos a setecientos: un cuarenta por ciento de sobrecosto. Es la segunda vez
que ocurre exactamente lo mismo, porque en 1911 la obra de 1910 había requerido una
ampliación del cuarenta y ocho por ciento.

Dos obras hidráulicas separadas por veinticinco años, la primera ejecutada por una comisión de vecinos que integraba el presidente municipal y la segunda por el municipio,
ambas con déficit y ambas cubiertas por decreto posterior. No es un dato sobre la
administración de nadie en particular: es lo que ocurre cuando se presupuesta una obra de
emergencia sobre un río que nadie midió.

En junio de 1936 aparece el cierre contable de esa obra: un decreto autoriza \textbf{treinta
y dos pesos con ochenta y siete centavos} a favor de la Dirección General de Obras Públicas
«a objeto de que pueda \textbf{terminar el pago} de los trabajos de defensas del pueblo de La
Caldera sobre el Río del mismo nombre», imputados a imprevistos del plan de obras.

Es un saldo, y por lo tanto una confirmación: \textbf{las defensas se construyeron}. No
quedaron en proyecto.

\subsection*{El informe técnico que nadie continuó}

Once meses después del espigón de emergencia, en diciembre de 1935, la Dirección General de
Obras Públicas elevó al gobernador un proyecto completo, con levantamiento de la zona,
planos y presupuesto. El decreto que lo aprueba transcribe el informe entero, y ese informe
es la mejor descripción del problema que este libro encontró en cualquier época.

\begin{ficha}{Informe sobre las obras de defensa del pueblo de La Caldera --- diciembre de 1935}
«El Río de La Caldera, a más o menos \textbf{cuatro kilómetros aguas arriba del pueblo}, se
divide en \textbf{tres brazos}, de los cuales el más amplio es el de la \textbf{margen
derecha}, siendo el más profundo y de mayor pendiente que los otros dos.

Este brazo \textbf{toma una vuelta brusca}, cambiando de rumbo después de pasar
inmediatamente junto al Pueblo, y recibe aguas abajo del mismo los otros dos brazos.

La ubicación en que se encuentra el pueblo hace que reciba en épocas de crecientes daños,
debido además a que \textbf{la mayor parte del caudal es transportado por el brazo de la
margen derecha}.

\ldots\ La solución consiste en la construcción de \textbf{defensas de piedra y ramas} algo
aguas arriba del punto donde se bifurcan los tres brazos, \textbf{con el objeto de modificar
la dirección de la corriente} de manera que el mayor caudal sea transportado por el brazo de
la margen izquierda. Complementando estas defensas, también se ha proyectado en la margen
izquierda \textbf{un canal profundo y angosto} a través de la punta que existe \ldots\ con
este se espera que por el efecto erosivo del agua desaparezca en breve tiempo dicha punta.

Con las obras descriptas se podrá librar al pueblo de posibles daños \textbf{y adquirir
experiencia sobre el comportamiento del Río, para que pasada la venidera época de crecientes
se puedan proyectar obras más duraderas}.»

Se autorizan \textbf{\$3.757,50}, por administración y con encuadre en el artículo 83 inciso
b de la Ley de Contabilidad --- ejecución directa por razones de urgencia ---, con cargo de
dar cuenta a la Legislatura.

\textit{El cuadro de presupuesto impreso no cierra: los dos ítems suman 3.757,50 y el total
bajo la raya dice 5.757,50. Verificado contra la imagen, el error es del original, publicado en 1936. La
cifra operativa, repetida en el informe y en los artículos 1º y 2º, es 3.757,50.}
\end{ficha}

Léase lo que ese informe entiende, porque es notable para 1935 y sigue siendo cierto.

\textbf{Identifica la causa exacta.} El río se bifurca en tres brazos cuatro kilómetros
aguas arriba; el de la margen derecha se lleva la mayor parte del caudal y toma una curva
brusca justo al pasar por el pueblo. El problema no es la crecida: es el reparto del caudal
entre brazos, decidido cuatro kilómetros antes.

\textbf{Propone actuar sobre la causa y no sobre el efecto.} No se trata de amurallar la
orilla junto a las casas sino de intervenir arriba, en la bifurcación, para mandar el agua
por el brazo izquierdo. Es --- en la lengua de 1935 --- una intervención sobre la dinámica de
la cuenca y no sobre el punto de daño.

\textbf{Y declara su propia provisionalidad con una honestidad que conviene subrayar.} Las
obras servirán, dice, para librar al pueblo de daños \textbf{y para adquirir experiencia
sobre el comportamiento del río}, de modo que después «se puedan proyectar obras más
duraderas».

Es el tercer acto de 1935 sobre la misma obra ---enero, febrero y diciembre--- y el segundo, después del de enero, en el que el Estado consigna por escrito que lo que está haciendo es provisorio y que la
solución definitiva requiere un estudio posterior.

\textbf{Ochenta y seis años después, la sentencia del amparo ordenará un plan de manejo que
privilegie infraestructura verde por sobre la gris y que atienda la cuenca y no el punto de
daño.} Es, con otras palabras y otro fundamento jurídico, lo que este informe proponía en
1935. La diferencia es que en 2021 lo ordenó un juez, y que en 1935 lo escribió el propio
organismo técnico del Estado y nadie lo continuó.

\subsection*{Un puente que no}

Cuatro meses antes, en agosto de 1935, la Dirección de Vialidad había tratado otro asunto del
departamento. El punto 11 de su acta dice: «\textbf{Vecinos de la Caldera} solicitan la
construcción de \textbf{un puente de unión entre dicho Pueblo y el Camino Nacional
Salta-Jujuy}».

La resolución, íntegra: contestar a los peticionantes que el Directorio ha tomado debida nota
y que dispondrá lo que corresponda \textbf{en su oportunidad}.

En la misma sesión, otro punto del acta consigna que las partidas de Vialidad están
sobreinvertidas y resuelve reducir el personal de las cuadrillas y gestionar un anticipo de
fondos. La respuesta dilatoria cae dentro de ese ajuste y no requiere otra explicación.

Pero el contraste queda registrado, y es el de este libro entero. \textbf{Los vecinos pidieron
un puente que los uniera al camino nacional y no lo obtuvieron. El Estado sí encontró
fondos, diez meses después, para terminar de pagar las defensas del río.} Una obra une al
pueblo con el afuera; la otra lo protege del agua. La segunda se hizo.

No es que una fuera más importante que la otra. Es que sólo una de las dos era, además,
protección de un camino nacional y de una traza que el capítulo \ref{cap:infraestructura} muestra que sirve a algo
más grande que el pueblo.

\textbf{Y el estudio que el decreto de enero de 1935 consideraba necesario para la solución
definitiva sigue sin aparecer noventa y un años después.} Lo que se hizo, siempre, fue el
espigón de emergencia --- y siempre costó más de lo previsto.

\section{1938: ahora piden los de la otra margen}

El informe técnico de 1935 proponía mandar el mayor caudal por el brazo de la margen
izquierda para librar al pueblo, que está sobre la derecha. Tres años después, quienes viven
sobre la margen izquierda piden que los defiendan.

\begin{ficha}{Decreto Nº 1502 del 13 de enero de 1938 --- Expediente 68-D}
Visto la nota suscrita por los \textbf{propietarios y vecinos del partido de La Calderilla},
jurisdicción del departamento de La Caldera, «en la cual solicitan al Gobierno de la
Provincia la construcción de \textbf{defensas sobre la margen izquierda del Río Caldera}»; y
considerando:

«Que el petitorio de referencia lo fundan en la urgente necesidad de que se construyan
defensas \ldots\ \textbf{en salvaguardia de las plantaciones pertenecientes a los
agricultores de esa jurisdicción}, máxime cuando \textbf{todos los años en la época de
lluvias} las aguas del Río Caldera amenazan las propiedades con eminente peligro de perjuicio
para los cultivos.»

Obra «Defensa Río Caldera», \textbf{cincuenta metros}, \$250, por administración, imputada a
la partida «Defensa en los Ríos Para Poblaciones» de la Ley 386. Firman Luis Patrón Costas y
Carlos Gómez Rincón.
\end{ficha}

Cuatro observaciones.

\textbf{La obra de 1935 desplazó el problema, no lo resolvió.} El informe técnico había
propuesto explícitamente modificar la dirección de la corriente «de manera que el mayor caudal
sea transportado por el brazo de la margen izquierda». Tres años más tarde, los propietarios
de esa margen piden defensas.

Este libro \textbf{no afirma} que una cosa haya causado la otra: el decreto de 1938 no
menciona las obras anteriores, los vecinos de La Calderilla dicen que el río los amenaza
«todos los años en la época de lluvias» --- es decir, no describen un problema nuevo --- y una
crecida perjudica ambas márgenes sin necesidad de que nadie la haya desviado. Establecerlo
exigiría los croquis de las dos obras, que el apéndice \ref{cap:pedidos} reclama.

Pero la secuencia queda registrada, y es la que este capítulo viene describiendo: se
interviene sobre un punto del cauce, aparece un reclamo en otro punto, se interviene ahí.

\textbf{Segunda: cambia quién pide y qué defiende.} En 1910 pidieron los vecinos del pueblo
y se trataba de proteger la población y el camino. En 1938 piden los propietarios de La
Calderilla y se trata de \textbf{salvaguardar plantaciones}. La defensa del río deja de ser
urbana y pasa a ser agrícola, en un departamento cuya superficie agropecuaria el capítulo \ref{cap:tierra}
verá desaparecer.

\textbf{Tercera: la escala se achica.} Dos mil pesos en 1910, tres mil setecientos en el
proyecto de 1935, quinientos en el espigón, doscientos cincuenta ahora por cincuenta metros
de defensa. Y \textbf{sigue siendo por administración y por urgencia}: el artículo que la
imputa se llama, textualmente, «Defensa en los Ríos Para Poblaciones».

\textbf{Y cuarta: existía una partida presupuestaria estable para esto.} La Ley 386 ---de 1936, registrada en el Digesto como Ley 1664--- tenía un ítem específico. No era una emergencia: era un renglón del presupuesto provincial.

\section{1940: el río se corrió solo}

El pedido de los vecinos de La Calderilla de 1938 tuvo continuación, y el decreto que la
ordena contiene el dato más revelador de toda esta serie.

\begin{ficha}{Decreto Nº 3.454 del 10 de febrero de 1940 --- Expediente 68-D/938}
La Dirección General de Obras Públicas solicita autorización «con carácter de urgencia» para
una defensa en el Río de la Caldera. Del informe de fojas 9 se desprende que \textbf{las obras
debieron ejecutarse el año anterior y se suspendieron porque el río desvió su curso}, y que
se estudiaba entonces la construcción de un puente.

Se autorizan \$500 para un \textbf{reparo provisorio} de defensa \textbf{en la margen
izquierda}, en el lugar denominado \textbf{La Calderilla}. Imputación: Ley 386, «Defensa En
Los Ríos Para Poblaciones». Firman Luis Patrón Costas y Carlos Gómez Rincón.
\end{ficha}

\textbf{El río desvió su curso y hubo que suspender la obra.} Es la primera vez que un acto
administrativo de esta serie reconoce que el cauce se mueve por sí mismo, y que esa
movilidad vuelve inútil una defensa proyectada un año antes.

Eso ilumina hacia atrás todo el capítulo. El informe técnico de 1935 había propuesto
\textbf{modificar deliberadamente la dirección de la corriente} para mandar el caudal al
brazo izquierdo. En 1938 los propietarios de esa margen pidieron defensas. En 1939 el río se
corrió solo. Y en 1940 se construye un reparo «provisorio» sobre esa misma margen.

Este libro \textbf{no puede establecer} si el desvío de 1939 fue efecto de las obras de 1935,
consecuencia natural de una crecida o ambas cosas. Lo que el expediente muestra es que
\textbf{el Estado intervino sobre un cauce móvil durante treinta años sin haberlo medido}, y que cuando el río se movió por su cuenta la respuesta fue, otra vez, un reparo
provisorio de quinientos pesos.

Conviene retener además la otra frase del informe: se estudiaba \textbf{la construcción de un
puente}. Cinco años después de que los vecinos pidieran uno a Vialidad y recibieran por
respuesta que se dispondría lo que correspondiera «en su oportunidad».

\section{1945: nadie se presentó}

Diez años después del espigón de emergencia, la Provincia intentó hacer las cosas como
corresponde: proyecto, presupuesto y licitación pública. No le salió.

\begin{ficha}{Licitación Pública Nº 1147 y Decreto Nº 9264-H, 1945}
El 25 de septiembre de 1945 la \textbf{Dirección General de Hidráulica} llama a licitación
pública «para las obras de defensa a efectuarse en el Río de La Caldera». Presupuesto oficial:
\textbf{\$2.515}. Apertura de propuestas: 10 de octubre, a las once, ante el Escribano de
Gobierno.

El 31 de octubre, el Decreto 9264-H deja constancia de que «en la licitación de referencia se
han llenado todos los requisitos exigidos por la Ley de Contabilidad y \ldots\ \textbf{a la
misma no se ha presentado ningún proponente}».

Y de que «por lo avanzado del año y la naturaleza del trabajo, el llamado a una nueva
licitación dilataría el trámite \textbf{haciendo imposible, seguramente, la ejecución de las
obras antes de que comience la época de lluvias}».

Se autoriza entonces la ejecución \textbf{por administración}, al amparo del artículo 83 de la
Ley de Contabilidad, por \textbf{\$2.904,83} --- con diez por ciento de inspección y cinco de
imprevistos ya incluidos. Imputación: Ley 712, Partida 14, «Obras de Defensa en los ríos para
poblaciones».
\end{ficha}

\textbf{Primera: no hubo oferentes.} Dos mil quinientos pesos de obra hidráulica a una veintena de kilómetros de la capital, y ninguna empresa se presentó. El aviso ofrecía los pliegos a cinco
pesos en la calle Mitre y publicaba el llamado en el Boletín con quince días de antelación.
El mercado de obra pública sencillamente no llegó.

\textbf{Segunda: el fundamento de la excepción es el calendario.} La contratación directa se
justifica porque una nueva licitación haría imposible ejecutar antes de las lluvias. Es el
mismo argumento que en 1910, en 1935 y en 1940: \textbf{la urgencia es la regla, no la
excepción}, y el río no espera al procedimiento.

\textbf{Tercera: y el costo volvió a subir antes de empezar.} El presupuesto licitado era de
\$2.515; el autorizado por administración, de \$2.904,83. Un quince por ciento más, esta vez
declarado de antemano --- son la inspección y los imprevistos ---, a diferencia de 1911 y 1935,
cuando el sobrecosto apareció después.

\subsection*{Lo que sí se votó en 1935}

Antes de seguir conviene poner al lado de esas cifras una ley del mismo año, porque cambia lo
que significan.

\begin{ficha}{Ley Nº 172 --- sancionada el 7 de junio de 1935, promulgada el 14}
«\textit{Art.\ 1º.--- Autorízase al Poder Ejecutivo para invertir la suma de \textbf{7.000,---}
en reformar y ampliar la \textbf{Iglesia Parroquial del Departamento de La Caldera}.}»

Las obras se ejecutarán según los planos y pliegos del ingeniero Antonio Mille; el gasto sale
de Rentas Generales. Sancionan C.\ Patrón Uriburu y Ricardo Solá; promulga Aráoz. B.O.\ Nº
1.589, página 46.
\end{ficha}

Puesta contra los otros tres números de 1935, la ley ordena la escena entera:

\begin{itemize}[leftmargin=*,itemsep=1pt]
\item \textbf{\$500} ---ampliados a \$700--- para el espigón de emergencia «pie de gallo» que impidiera que el río entrara a la calle principal.
\item \textbf{\$3.757,50} para el proyecto completo de Obras Públicas, con levantamiento, planos y presupuesto, que identificaba la bifurcación y proponía intervenir ahí. El Boletín no registra su ejecución.
\item \textbf{\$7.000} para reformar y ampliar la iglesia.
\end{itemize}

\textbf{La Legislatura votó para la iglesia diez veces lo que el Ejecutivo liquidó para
defender el pueblo del río, y casi el doble de lo que costaría la solución que su propia
repartición técnica proyectó en diciembre.} Los tres actos son del mismo año y del mismo Boletín.

\begin{cautela}
Esto no dice que la iglesia no hiciera falta ni que las dos decisiones fueran alternativas: no
salían de la misma partida ---la iglesia, de Rentas Generales por ley; las defensas, de
las vías de urgencia de la Ley de Contabilidad por decreto--- y las tomaron órganos distintos. Lo que
dice es que \textbf{en 1935 había siete mil pesos disponibles para obra en La Caldera}, y que
el criterio por el cual el río recibió setecientos ---más los \$3.757,50 del proyecto que se autorizó en diciembre--- no está escrito en ninguna parte.
\end{cautela}

\subsection*{Un año antes, una obra chica, y todo el procedimiento a la vista}

Vale poner al lado de esa licitación desierta una obra del mismo período que sí completó el
ciclo. No es de defensa ni toca el río: es la refección del local de la Comisaría de La
Caldera, y sirve como patrón contra el cual leer todo lo demás.

{\small
\begin{longtable}{@{}p{2.5cm}p{2.0cm}p{7.4cm}@{}}
\toprule
\textbf{Acto} & \textbf{Fecha} & \textbf{Contenido} \\
\midrule
\endhead
Dto.\ 2424 & 04/03/1944 & Dispone la licitación pública, junto con la comisaría de Palomitas (Campo Santo) \\
Acta de licitación & 11/04/1944 & \textbf{Un solo oferente}: el contratista Francisco Crescini \\
Dto.\ 2965-H & 27/04/1944 & Adjudica. La Caldera, \textbf{\$5.018,02} ---15\,\% sobre el presupuesto oficial---; Palomitas, \$4.187,01 ---19\,\%---. Total \$9.205,03, Ley 712, partida «Reparaciones de Comisarías de Campaña» \\
Dto.\ 2969-H & 27/04/1944 & \$40 al diario \textit{La Provincia} por el aviso \\
Dto.\ 2974-H, art.\ 3º & 27/04/1944 & \$40 al diario \textit{Norte} por el mismo aviso \\
Dto.\ 4417-H & 04/09/1944 & Certificado Nº 1 \textbf{final}, \$5.018,02, con retención del 10\,\% en garantía \\
Dto.\ 4681-H & 27/09/1944 & Devuelve el depósito en garantía de las dos obras, \$460,25: el 5\,\% exacto de los \$9.205,03 adjudicados. Es otro concepto que la retención del 10\,\% sobre el certificado \\
Dto.\ 4727-H & 29/09/1944 & Certificado \textbf{adicional} Nº 1 (final), \textbf{\$555,78} \\
\bottomrule
\end{longtable}}

Tres cosas se leen ahí, y las tres reaparecen en las obras del río.

\textbf{El sobrecosto está en dos tramos.} Quince por ciento en la adjudicación, y un adicional
de \$555,78 sobre un certificado ya rotulado «final» ---un once por ciento más---. Sumados sobre
el presupuesto oficial, el sobrecosto total ronda el veintiocho por ciento. Puesto en la serie:
cuarenta y ocho por ciento en 1911, cuarenta en 1935, quince y medio en 1945.

\textbf{No hubo competencia, y el decreto lo dice.} Un solo oferente, y la adjudicación se funda
en que las ofertas «son equitativas dada la carestía de los materiales». Es el argumento de la
inflación de guerra, enunciado sin rodeos.

\textbf{Y el orden de los actos no cierra.} La devolución del depósito en garantía se autoriza
el 27 de septiembre; el certificado adicional final, el 29. \textbf{Se devolvió la garantía dos
días antes de certificar el último trabajo.}

\begin{cautela}
El presupuesto oficial de cada obra \textbf{no está en el Boletín}. Se puede despejar de los
porcentajes ---\$4.363,50 para La Caldera y \$3.518,50 para Palomitas--- y así se hizo aquí, pero
es un cálculo de este libro y no un dato publicado. El decreto tampoco dice en qué estado
estaba el edificio, ni si el local era propio o alquilado. El expediente más antiguo del legajo
es de 1942: la refacción venía tramitándose desde hacía al menos dos años.
\end{cautela}

\textbf{Una cuarta observación sobre la licitación de 1945: el aviso no dice dónde.} «Las obras de defensa a efectuarse en el Río de La
Caldera», sin tramo, sin tipo de obra y sin indicar qué población se protege. La
caracterización aparece recién en la imputación presupuestaria. Es la misma opacidad de
criterio que el capítulo \ref{cap:opacidad} describe para el presente: el acto se publica, el contenido
operativo no.

\section{1951: la defensa se vuelve permanente}

En febrero de 1951 la serie del río cambia de escala y de organismo, y esta vez el cambio
queda escrito en la imputación presupuestaria.

\begin{ficha}{Decretos Nº 5268-E y 5269-E, 2 de febrero de 1951}
El \textbf{5268-E} aprueba el proyecto y presupuesto elaborado por el Departamento de
Ingeniería de la \textbf{Administración General de Aguas de Salta} para la obra «Defensas
sobre el río de La Caldera», por \textbf{\$27.781,95}, con el cinco por ciento de imprevistos
incluido.

El \textbf{5269-E}, dictado el mismo día, autoriza \textbf{tres defensas de rama y piedra sobre
la margen izquierda del río Vaqueros}, según croquis de la Inspección destacada, hasta un
máximo de \$10.000 y por vía administrativa.

\textbf{Ambos se imputan a la misma partida:} Anexo I, Inciso IV, Principal 1, Parcial b),
«\textbf{Obras de defensa permanente}», Partida 7, «en Rosario de Lerma, \textbf{La Caldera},
Capital, Cerrillos y Campo Santo».
\end{ficha}

\textbf{Primera: el monto se multiplicó por diez.} Quinientos pesos el espigón de 1935,
doscientos cincuenta las defensas de La Calderilla en 1938, dos mil novecientos la obra de
1945. Veintisiete mil ochocientos en 1951. Aun descontando la inflación del período, es otra
categoría de intervención.

\textbf{Segunda: aparece un organismo técnico especializado.} No es ya la Dirección General de
Obras Públicas ni Hidráulica: es la \textbf{Administración General de Aguas de Salta}, con su
propio Departamento de Ingeniería, su Honorable Consejo y sus resoluciones previas. El Estado
provincial se dotó de un aparato para el agua.

\textbf{Tercera, y es la que importa: la defensa dejó de ser una emergencia.} La partida se
llama «obras de defensa \textbf{permanente}» y nombra a cinco departamentos, entre ellos La
Caldera. En 1938 el rubro presupuestario se llamaba «Defensa en los Ríos Para Poblaciones» y
financiaba reparos provisorios; trece años después hay un renglón fijo para obra permanente.

\textbf{Y cuarta: siguen siendo defensas de rama y piedra.} El decreto lo dice: tres defensas
«de rama y piedra» sobre la margen izquierda del Vaqueros, según croquis de la inspección. La
partida es permanente; la técnica, la misma de 1910.

\begin{cautela}
El estudio de cuenca que el informe de 1935 consideraba necesario para «proyectar obras más
duraderas» no aparece tampoco aquí. Lo que estos decretos aprueban es un proyecto de obra
--- planos y presupuesto de tres defensas ---, no un estudio del comportamiento del río.

\textbf{Es un salto institucional y presupuestario, no de conocimiento.} Y conviene tenerlo
presente al leer el capítulo \ref{cap:amparo}: en 2021 una jueza vuelve a ordenar lo que en 1935 el propio
Estado había escrito que hacía falta.
\end{cautela}

\section{1980: hormigón, y presupuesto nacional}

Entre 1951 y 2001 esta serie tiene un hueco de medio siglo que el relevamiento no cubre. Un
aviso de abril de 1980 permite ver qué había pasado en el medio.

\begin{ficha}{Licitación Pública Nº 9/80 --- Dirección Nacional de Vialidad, 5º Distrito}
Obra: \textbf{Ruta 9, tramo Río Caldera --- Alto de la Sierra, construcción de defensas de
hormigón}, en jurisdicción de la provincia de Salta.

Presupuesto oficial: \textbf{\$62.264.460}. Precio del pliego: \$12.500. Plazo de obra:
\textbf{cuatro meses}. Apertura: 9 de mayo de 1980, en la sede del 5º Distrito, Pellegrini
715. Firma: Néstor Enrique Berto, Jefe de la División Licitaciones y Contratos.

El aviso se publicó cinco días seguidos, del 22 al 28 de abril.
\end{ficha}

\textbf{Primera: el material.} Rama y piedra en 1910, en 1938 y todavía en 1951. En 1980,
\textbf{hormigón}. Es la primera obra de defensa de la serie que no se hace con lo que hay en
el lugar.

\textbf{Segunda: quien la paga.} No es la Provincia ni el municipio: es la \textbf{Dirección
Nacional de Vialidad}, con presupuesto nacional y licitación pública en regla. Cuatro meses
de plazo y pliego a doce mil quinientos pesos.

\textbf{Y tercera, que es la que importa: lo que se defiende no es el pueblo.} La obra es
sobre la \textbf{Ruta 9}, en el tramo que va del río Caldera al Alto de la Sierra. Se protege
la traza vial que une Salta con Jujuy.

Póngase esa cifra al lado de las otras de esta serie. En 1935, quinientos pesos para el
espigón que defendía las casas del pueblo. En 1938, doscientos cincuenta para cincuenta
metros de defensa de las plantaciones de La Calderilla. En 1945, dos mil novecientos, y
nadie se presentó a la licitación. En 1951, veintisiete mil ochocientos, ya con partida
permanente.

Las monedas son distintas y la comparación directa no procede. Pero \textbf{lo que sí procede
es la comparación de quién pone el dinero}: cuando lo que está en riesgo es el camino
nacional, aparece un organismo nacional con licitación pública, cuatro meses de plazo y
hormigón. Cuando lo que está en riesgo son las casas, la obra se hace por administración, con
rama y piedra, y con el argumento de que hay que llegar antes de las lluvias.

Es la misma asimetría que el capítulo \ref{cap:infraestructura} describe, y aquí aparece en una sola línea
presupuestaria.

\subsection*{Tres cuencas en tres semanas}

Cuatro meses después de la licitación de Vialidad Nacional, entre el 27 de agosto y el 19 de
septiembre de 1980, la Provincia aprobó tres obras de encauzamiento en el departamento. Las
tres por la Administración General de Aguas, las tres \textbf{por vía administrativa}.

\needspace{7\baselineskip}
{\sffamily\bfseries\footnotesize Obras de encauzamiento aprobadas en el departamento, agosto--septiembre de 1980\par}
\vspace{0.5em}
{\small
\setlength{\LTpre}{0.6em}\setlength{\LTpost}{1.2em}
\begin{longtable}{@{}p{2.3cm}p{5.4cm}p{4.2cm}@{}}
\toprule
\textbf{Resolución} & \textbf{Obra} & \textbf{Presupuesto oficial} \\
\midrule
\endhead
880-D, 27 ago. & Encauzamiento sobre el \textbf{río Wierna} & \$335.110.551 \\
914, 3 set. & Encauzamiento y defensa sobre el \textbf{río La Caldera}, zona \textbf{Campo Alegre} & \$107.291.751 \\
943-D, 19 set. & Encauzamiento sobre el \textbf{río Vaqueros} & \$85.219.679 \\
\midrule
\multicolumn{2}{@{}l}{\textbf{Total}} & \textbf{\$527.621.981} \\
\bottomrule
\end{longtable}
}

\textbf{Primera: son las tres cuencas del departamento a la vez.} El Wierna, el La Caldera y
el Vaqueros. En veintitrés días, el Estado provincial decidió intervenir sobre los tres
cursos principales del territorio.

\textbf{Segunda: la del río La Caldera se hace en la zona de Campo Alegre.} Es decir, sobre el
área del embalse expropiado en 1972, que ese mismo año recibía trescientos millones para
proyectar su planta potabilizadora. \textbf{Se protege el sistema que abastece a la ciudad de
Salta.}

\textbf{Y tercera: ninguna se licitó.} Las tres resoluciones autorizan a la A.G.A.S.\ a
ejecutar «por vía administrativa», sin llamado público. Quinientos veintisiete millones de
pesos en obra hidráulica contratada directamente.

El contraste con la obra de la Ruta 9 es exacto y vale enunciarlo. \textbf{Vialidad Nacional
licitó públicamente las defensas del camino; la Provincia contrató directamente las de los
ríos.} Y aunque las cifras no son comparables entre sí sin corregir por inflación, sí lo es el
procedimiento: el mismo año, el mismo territorio, dos estándares distintos.

\subsection*{Y había alguien midiendo la lluvia}

Este capítulo repite, a lo largo de ciento dieciséis años, que el Estado intervino sobre estos ríos sin medirlos. \textbf{El capítulo \ref{cap:ribera} corrige esa afirmación}: hubo cinco estaciones de aforo en la cuenca, y sus registros están publicados. Lo que este capítulo documenta no es la ausencia del dato sino \textbf{que ninguna de las obras que siguen lo invoca}. Un
decreto de junio de 1980 obliga a precisar esa afirmación.

\begin{ficha}{Decreto Nº 876 del 27 de junio de 1980}
Expediente C/19-02659/80. El Ministerio de Economía autoriza el pago por «trabajos de
observaciones meteorológicas y pluviométricas durante el año 1980» a los observadores de una
red provincial.

Entre ellos, el del \textbf{Observatorio Pluviométrico Los Yacones} --- paraje del
departamento La Caldera ---: \textbf{Daniel Maidana}, por sesenta mil pesos
mensuales durante doce meses.

El decreto abona en total \$15.480.000 a toda la red de observadores.
\end{ficha}

\textbf{Había una estación pluviométrica en el departamento, con un observador pago por la
Provincia, al menos en 1980.} Alguien anotó cuánto llovía en Los Yacones, todos los días de
ese año, y cobró por hacerlo.

Eso obliga a distinguir dos cosas que este libro venía diciendo juntas.

\textbf{Lo que importa no es la medición de la lluvia sino el aforo del río.} Son datos
distintos. La pluviometría registra cuánta agua cae; el aforo, cuánta lleva el cauce. El
artículo 8 del Decreto 1.989/02 exige lo segundo, con series de cincuenta años, para
determinar una línea de ribera. El decreto de 1924 declaraba no tener aforos, no declaraba no
tener pluviómetros.

Pero la distinción juega en contra del Estado, no a su favor. \textbf{Si en 1980 había una
red provincial de observadores con presupuesto y nómina, entonces la capacidad institucional
de medir existía.} Lo que no existió fue la decisión provincial de aplicarla al caudal de estos cauces: las estaciones de aforo de la cuenca eran nacionales, y en 1980 la de El Angosto atravesaba uno de sus tramos sin registro (capítulo \ref{cap:ribera}).

Y ahora sabemos cuánto duró. El Plan de Manejo de 2024 registra a \textbf{Los Yacones}
entre las estaciones pluviométricas de la cuenca, con serie de \textbf{1972 a 1990}, cota
1.550 metros, titular la Dirección General de Aguas de Salta, \textbf{estado: cerrada}.

De modo que el observador al que la Provincia le pagaba sesenta mil pesos mensuales en 1980
trabajaba en una estación que funcionó dieciocho años y cerró en 1990. \textbf{Y era la
única estación de origen provincial de toda la cuenca}: las cinco de aforo y al menos tres de las otras cuatro pluviométricas pertenecían a organismos nacionales.

\section{2001: un arroyo que se llama Urquiza}

Después de 1980, el siguiente eslabón documentado de esta serie es de noviembre de
2001, y trae un dato que este libro no buscaba.

\begin{ficha}{Resolución SOP Nº 516, Secretaría de Obras Públicas}
Expediente 34-2.218/01. Aprueba la documentación técnica preparada por la Unidad de
Infraestructura de Desarrollo para la obra «\textbf{Encauzamiento Arroyo Urquiza\index{Urquiza, Francisco S.}, Vaqueros,
departamento La Caldera}», con un presupuesto oficial de \textbf{\$7.360}.

Con encuadre en el artículo 12 de la Ley 6.838, adjudica a la Empresa \textbf{Irac S.R.L.}
por \$7.356,80 con IVA incluido. Plazo: \textbf{diez días corridos}.
\end{ficha}

\textbf{Hay un arroyo del departamento que lleva el apellido Urquiza\index{Urquiza, Francisco S.}.}

Este libro no puede establecer si el nombre viene de Francisco S.\ Urquiza\index{Urquiza, Francisco S.} --- el que en 1921
hizo deslindar Cerro de Buena Vista, en 1924 obtuvo cien litros por segundo del río Wierna,
en 1928 fue detenido por hurto de agua, en 1931 le ganó el juicio a la Provincia y en 1944
figuraba como el mayor propietario del departamento --- o de un homónimo anterior. El
capítulo \ref{cap:fincas} documenta que en este territorio los topónimos salen de los apellidos de quienes
tuvieron la tierra: Vaqueros de un poseedor colonial, Wierna de un litigante de 1908. El
patrón está establecido; la atribución concreta, no.

Lo que sí puede decirse es que \textbf{setenta años después del pleito por el agua, el Estado
provincial encauza un arroyo que lleva ese nombre}, con documentación técnica propia,
licitación encuadrada y diez días de plazo.

Y dos observaciones sobre la obra misma.

\textbf{La escala es minúscula.} Siete mil trescientos sesenta pesos, diez días corridos. El
capítulo documenta obras de veintisiete mil pesos en 1951; medio siglo más tarde, en moneda
muy distinta, se encauza un arroyo por el valor de un automóvil usado.

\textbf{Y la adjudicación va por el artículo 12 de la Ley 6.838}, que es la vía abreviada de
contratación. Como en 1935, como en 1940, como en 1945: \textbf{la excepción al procedimiento
ordinario es la forma normal de intervenir sobre estos cauces.}

\section{2014 y 2015: se sigue moviendo suelo en el cauce}

\begin{ficha}{Contratación directa --- B.O.\ Nº 19.523, 21 de abril de 2015}
La Compañía Salteña de Agua y Saneamiento contrata directamente, por el artículo 13 h de la Ley 6838, el \textbf{``servicio de equipo pesado tipo 320 para movimiento de suelo en río La Caldera''}, con destino a las \textbf{localidades de La Caldera y La Calderilla}. Fecha de contratación: 10 de abril de 2015. Proveedor: \textbf{Valle Vial S.A.} Monto: \$47.100 más IVA. Expediente 7766/15.
\end{ficha}

\begin{cautela}
El objeto es literal y conviene leerlo así: \textbf{movimiento de suelo en el río}. Una excavadora de la clase que la industria designa como ``tipo 320'' trabajando en el cauce del río La Caldera, contratada por la empresa de agua y saneamiento, tres años y ocho meses antes del primer aluvión.

No hay nada anómalo en que una empresa de agua intervenga sobre el lecho de un río del que capta o cerca del cual tiene infraestructura: las tomas se embancan, los conductos se descubren, las crecidas depositan material. Y este libro no sabe qué se movió, hacia dónde ni con qué autorización.

Lo que registra es que la intervención mecánica sobre el cauce del río La Caldera \textbf{es una constante documentada del período}, no un episodio de 2026. En 2014 la Provincia gastó casi dos millones en proteger de las crecidas las cámaras de captación; en 2015 la empresa de agua movió suelo en el lecho con equipo pesado; el municipio construía badenes y gaviones por contrato de Vialidad desde 2014; y en febrero de 2026 informaba bateas, encauzamiento y muros de contención con maquinaria pesada.

Doce años de intervención sobre el mismo cauce, por distintos organismos, sin que este libro haya encontrado un solo acto que evalúe el efecto acumulado de todas ellas. El artículo 171 del Código de Aguas exige esa evaluación previa para los permisos de extracción de áridos. Para las obras del propio Estado sobre el cauce, no se encontró norma equivalente ni constancia de que se haya hecho.
\end{cautela}

\subsection*{Proteger la toma de agua del río que la alimenta}

En marzo de 2014, cuatro años y nueve meses antes del primer aluvión, la Provincia adjudicó una obra sobre el mismo río.

\begin{ficha}{Resolución SOP Nº 114 del 11 de marzo de 2014 --- B.O.\ Nº 19.274}
Obra: \textbf{``Protección y Defensas de Cámaras de Inspección --- Sistema de Captación de Agua Acueducto Norte Río La Caldera''}, departamento La Caldera. Expediente 125-217.210/13.

Presupuesto oficial \$2.447.148,48, modalidad ajuste alzado. Procedimiento selectivo del 26 de diciembre de 2013, a cargo de la Dirección de Obras Hídricas y Saneamiento.

Ofertas admisibles: Dal Borgo Construcciones S.R.L., \textbf{IN.CO.VI.\ S.R.L.}, C.E.\ y BA.\ S.R.L., Eco Suelo S.A.\ y Romero Igarzabal S.R.L.

Adjudicada a \textbf{Dal Borgo Construcciones S.R.L.} en \$1.824.587,97 --- un 25,44 por ciento por debajo del presupuesto ---, con plazo de sesenta días corridos.
\end{ficha}

\begin{cautela}
El objeto de la obra dice, en lenguaje administrativo, algo que este libro viene diciendo de otro modo: \textbf{había que proteger de la crecida la toma de agua del acueducto}.

Las cámaras de inspección del sistema de captación estaban expuestas al río, y en diciembre de 2013 la Provincia decidió que necesitaban defensas. Es el Acueducto Norte, que en 2014 captaba del río ---según su prestadora, a la altura de la curva de El Gallinato--- y conducía el agua hacia la cisterna de El Huaico; desde 2021 el mismo nombre designa también el sistema nuevo que toma del dique Campo Alegre (capítulo \ref{cap:expropiacion}), y el mismo curso sobre el cual, en ese mismo período, se determinaban líneas de ribera en ocho inmuebles.

De modo que en 2013 y 2014 el Estado provincial hizo tres cosas a la vez sobre el río La Caldera: relevó sistemáticamente sus márgenes, gastó casi dos millones de pesos en proteger su infraestructura de agua de las crecidas, y --- a través de otras dependencias --- siguió tramitando la aprobación de fraccionamientos sobre sus afluentes.

Este libro no sostiene que esas tres líneas de acción se contradijeran ni que alguien debiera haberlas coordinado y no lo hizo. Registra que \textbf{ocurrieron simultáneamente, en el mismo territorio, bajo el mismo gobierno}, y que la primera y la segunda demuestran que el riesgo hídrico del río La Caldera no era una hipótesis para la administración: era un problema presupuestado.
\end{cautela}

\begin{cautela}
Un dato que corresponde consignar sin sacar conclusión. Entre las cinco empresas cuyas ofertas fueron declaradas admisibles figura \textbf{IN.CO.VI.\ S.R.L.} Es la misma sociedad que, en ese mismo trimestre, aparece en el edicto de renovación de la concesión de la cantera de áridos ``El Vaquero'', nueve hectáreas de \textbf{tierra fiscal} en Vaqueros, según se documenta en el capítulo \ref{cap:resistencias}.

No hay nada irregular en que una empresa de construcción vial e hídrica explote además una cantera de áridos: es el insumo natural de su actividad, y ambas cosas se obtienen por procedimientos públicos distintos e independientes. Lo que este libro registra es la \textbf{concurrencia} de esa firma en circuitos que se presentan como separados, porque es verificable y porque el capítulo \ref{cap:opacidad} sostiene que en este departamento los mismos nombres reaparecen.

Y no son dos apariciones sino \textbf{tres}. En diciembre de 2014 la Compañía Salteña de Agua y Saneamiento contrató directamente a \textbf{Incovi S.R.L.}, por el artículo 13 h de la Ley 6838, el ``alquiler y acarreo de reparación de Acueducto Norte'' con destino a la \textbf{localidad de Vaqueros}, por \$11.575 más IVA, expediente 6903/14. De modo que en el mismo año la misma sociedad figura como concesionaria de nueve hectáreas de áridos sobre tierra fiscal, como oferente admisible en la obra de protección del acueducto y como contratista directa de su reparación.

Establecer si hubo o no relación entre ambas posiciones requiere los dos expedientes, y este libro no los tiene.
\end{cautela}

\section{Septiembre de 2026: el río volvió a moverse}

Este capítulo se cerró tres veces y hubo que reabrirlo. La última, en la semana en que se cerraba este relevamiento.

\begin{ficha}{Resolución Nº 479 de la Secretaría de Obras Públicas, 3 de septiembre de 2026 (B.O.\ Nº 22.266, 10 de septiembre)}
El Gobierno provincial \textbf{da de baja el contrato} para construir defensas sobre el río La
Caldera, adjudicado por \textbf{\$216.819.417,19}. El convenio, firmado entre la Provincia y
la Municipalidad, fijaba un plazo de ejecución de \textbf{noventa días}.

Según la documentación oficial, \textbf{los cambios registrados en los niveles y en el cauce
del río hicieron que la solución técnica prevista ya no reuniera las condiciones necesarias}
para resolver el problema original. La resolución aprueba el acuerdo de disolución.

No se precisó cuál será el nuevo proyecto técnico.
\end{ficha}

Póngase esta resolución al lado del decreto de 1940 que este capítulo transcribe más arriba. Aquel dejó constancia de que unas obras de defensa proyectadas el año anterior se habían suspendido \textbf{porque el río había desviado su curso}.

\textbf{Ochenta y seis años separan los dos actos, y dicen lo mismo.} Un proyecto aprobado, un río que se mueve antes de que la obra empiece, y un acto administrativo que registra que la solución prevista ya no sirve.

Tres observaciones.

\textbf{Primera: ocurre siete meses después de la aprobación judicial del Plan de Manejo.} La
jueza aprobó el Plan el 12 de febrero de 2026 y fijó cuarenta y cinco días para que la
Provincia presentara la planificación de obras prioritarias. En septiembre, una de esas obras
\textbf{se da de baja por inadecuación técnica sobreviniente}.

\textbf{Segunda: el fundamento es exactamente el que la norma prevé y la serie venía mostrando.} El artículo 7 del
Decreto 1.989/02 manda considerar la «migración lateral por erosión natural» del cauce; el
informe técnico de 1935 identificó que el río se bifurca y se corre; el decreto de 1940 lo
constató en los hechos. \textbf{Que este río cambia de lugar está documentado desde hace
noventa años, y sin embargo cada proyecto se diseña como si el cauce fuera fijo.}

\textbf{Y tercera: la obra tiene una historia que el índice del Boletín deja ver, aunque no la cierra.} Los títulos de las resoluciones de la Secretaría de Obras Públicas publicadas desde 2023 ordenan la secuencia:

\begin{itemize}[leftmargin=*,itemsep=1pt]
\item \textbf{Resolución 896/23} (B.O.\ Nº 21.617, 22 de diciembre de 2023): aprueba el legajo técnico de la «Defensa río La Caldera (zona Bº Santiago Apóstol sección 2 y zona Nogalar 4)».
\item \textbf{Resolución 314/24}, del 19 de marzo de 2024 (B.O.\ Nº 21.681): aprueba el convenio de rescisión de esa obra con la Municipalidad. El mismo día, la \textbf{315/24} rescinde también la «Ampliación de red de agua barrio Santiago Apóstol».
\item \textbf{Resolución 951/24}, del 1 de octubre de 2024 (B.O.\ Nº 21.812): aprueba el legajo técnico de una nueva «Defensa sobre río La Caldera».
\item \textbf{Resolución 300/26}, del 4 de junio de 2026 (B.O.\ Nº 22.205): aprueba el legajo técnico de un «Encauzamiento en río La Caldera --- Municipio La Caldera».
\item \textbf{Resolución 479/26}, del 3 de septiembre de 2026: rescinde la obra «Defensas sobre río La Caldera --- La Caldera».
\end{itemize}

Encaja con lo que el intendente y el ministro de Infraestructura explicaron en julio de 2025 (capítulo \ref{cap:amparo}): un convenio de defensas de 2023, de 63 millones, rescindido por desactualizado ---la rescisión es de marzo de 2024--- para firmar otro a valores corrientes. \textbf{La obra dada de baja en 2026 es la de la Resolución 951/24.} El texto lo confirma: presupuesto oficial de \$216.819.417,19 ---el mismo monto que la 479/26 rescinde---, adjudicada por contratación abreviada a la Municipalidad de La Caldera, noventa días de plazo; y el propio Boletín vincula las dos resoluciones en su ficha. Y confirma algo más: \textbf{es la misma obra de 2023, rehecha a valores de agosto de 2024}. Ciento ochenta metros lineales de gaviones y colchonetas en las mismas dos zonas ---sesenta en el barrio Santiago Apóstol, sección~2, y ciento veinte en El Nogalar~4---, a pedido del intendente, con documentación revisada por Recursos Hídricos, «para mitigar desbordes y a la vez dificultar el avance del río La Caldera en dirección de los Barrios Santiago Apóstol y El Nogalar». Es El Nogalar~4, el loteo que el intendente describiría en enero de 2026 como construido dentro del lecho del río (capítulo \ref{cap:amparo}). Lo que la secuencia muestra, además, es la forma: \textbf{en treinta y tres meses, dos legajos de defensa rescindidos y un tercer proyecto aprobado tres meses antes de la segunda rescisión}. No se sabe si alguno de ellos es una de las medidas estructurales del Plan de Manejo ---la defensa del río La Caldera está presupuestada allí en trece mil seiscientos noventa y seis millones, sesenta y tres veces los 216,8 millones de la obra rescindida---.

\section{Lo que la serie deja}

Ciento dieciséis años de intervenciones sobre el mismo cauce: reparos en 1910, sobrecosto en
1911, espigón de emergencia y proyecto técnico en 1935, saldo en 1936, defensas sobre la otra
margen en 1938, reparo provisorio tras el desvío del cauce en 1940, licitación desierta y ejecución por administración en 1945, partida de defensa permanente y A.G.A.S.\ en 1951, defensas de hormigón de Vialidad Nacional sobre la Ruta 9 y tres encauzamientos provinciales por vía administrativa en 1980, encauzamiento del arroyo Urquiza\index{Urquiza, Francisco S.} en 2001, protección de las cámaras de captación en 2014, movimiento de suelo con
equipo pesado en 2015, bateas, encauzamiento y muros de contención en febrero de 2026, y la baja de un contrato de defensas en septiembre de ese año porque el cauce se había corrido.

Este libro no concluye de ahí que las obras hayan sido inútiles ni mal hechas: no tiene los
elementos para juzgarlo, y buena parte de ellas probablemente evitó daños que no podemos
contar.

Lo que sí puede decir es que \textbf{durante todo ese siglo la respuesta fue contener el
agua}, y que en ninguno de los actos relevados aparece la otra pregunta --- qué se construyó,
dónde y con qué autorización sobre las márgenes que había que defender.

Esa pregunta la formula por primera vez la sentencia de 2021. Y hay una coincidencia que
conviene tener presente al leer los dos capítulos que siguen: lo que esa sentencia ordena
--- atender la cuenca y no el punto de daño, privilegiar infraestructura verde sobre la gris
--- es, en otras palabras, lo que el informe técnico de diciembre de 1935 ya proponía.
